\section{Additional results on multiplicative domains and order}

\subsection{Abstract Schwarz maps and their multiplicative domains}

\begin{definition}[Schwarz inequality for a linear map]
    \label{def:abstract_schwarz_inequality}
    \lean{IsSchwarzMap}
    \leanok
    A complex-linear map $E:\MN{D}\to\MN{D}$ satisfies the Schwarz
    inequality if
    \[
        E(A^\dagger A)-E(A^\dagger)E(A)\succeq 0
    \]
    for every $A\in\MN{D}$.  This is
    \cite[Equation~(5.2)]{Wolf2012Quantum}.
\end{definition}

\begin{definition}[Abstract one-sided multiplicative domains]
    \label{def:abstract_one_sided_multiplicative_domains}
    \lean{SchwarzMap.rightMultiplicativeDomain,
        SchwarzMap.leftMultiplicativeDomain,
        SchwarzMap.multiplicativeDomain}
    \leanok
    For a complex-linear map $E:\MN{D}\to\MN{D}$, set
    \[
        \mc{A}_R(E)
          = \{A : E(XA)=E(X)E(A)\text{ for every }X\},
        \qquad
        \mc{A}_L(E)
          = \{A : E(AX)=E(A)E(X)\text{ for every }X\},
    \]
    and $\mc{A}(E)=\mc{A}_R(E)\cap\mc{A}_L(E)$.
    These are \cite[Equations~(5.11)--(5.12)]{Wolf2012Quantum}, with the
    source's convention that the subscript records the side on which $A$
    multiplies the varying matrix.
\end{definition}

\begin{theorem}[Abstract multiplicative-domain characterization]
    \label{thm:abstract_md_characterization}
    \lean{SchwarzMap.mem_rightMultiplicativeDomain_iff,
        SchwarzMap.mem_leftMultiplicativeDomain_iff,
        SchwarzMap.mem_multiplicativeDomain_iff}
    \leanok
    \uses{def:abstract_schwarz_inequality,
        def:abstract_one_sided_multiplicative_domains}
    If $E$ satisfies the Schwarz inequality, then
    \[
        A\in\mc{A}_R(E)
          \quad\Longleftrightarrow\quad
        E(A^\dagger A)=E(A^\dagger)E(A),
    \]
    and
    \[
        A\in\mc{A}_L(E)
          \quad\Longleftrightarrow\quad
        E(AA^\dagger)=E(A)E(A^\dagger).
    \]
    Consequently, $A\in\mc{A}(E)$ exactly when both equalities hold.
    This is \cite[Equations~(5.13)--(5.14)]{Wolf2012Quantum}.
\end{theorem}

\begin{proof}
    \leanok
    \uses{def:abstract_schwarz_inequality}
    Suppose first that equality holds at $A^\dagger A$.  Apply the Schwarz
    inequality to $A+tY$, first for real $t$ and then for purely imaginary
    $t$.  Set
    \[
        L=E(A^\dagger Y)-E(A^\dagger)E(Y),\qquad
        R=E(Y^\dagger A)-E(Y^\dagger)E(A),\qquad
        Q=E(Y^\dagger Y)-E(Y^\dagger)E(Y)\succeq0.
    \]
    For every real $t$, expansion and the equality at $A^\dagger A$ give
    \[
        t(L+R)+t^2Q\succeq0,
        \qquad
        t\,i(L-R)+t^2Q\succeq0.
    \]
    Nonnegativity for both signs of every real parameter forces
    \[
        L+R=0,\qquad L-R=0.
    \]
    Hence $L=R=0$, which is precisely
    \[
        E(A^\dagger Y)=E(A^\dagger)E(Y),
        \qquad
        E(Y^\dagger A)=E(Y^\dagger)E(A).
    \]
    Replacing $Y$ by $X^\dagger$ yields $E(XA)=E(X)E(A)$.
    The second characterization follows by applying the same argument to
    $A^\dagger$.  The converse implications follow by evaluating the defining
    product identities at $A^\dagger$.
\end{proof}

\begin{theorem}[Kraus maps satisfy the abstract Schwarz inequality]
    \label{thm:kraus_abstract_schwarz}
    \lean{KadisonSchwarz.krausMap_isSchwarzMap}
    \leanok
    \uses{def:kraus_map, def:unital_kraus,
        def:abstract_schwarz_inequality}
    Every unital Kraus map satisfies the abstract Schwarz inequality of
    Definition~\ref{def:abstract_schwarz_inequality}.
\end{theorem}

\begin{proof}
    \leanok
    \uses{thm:kadison_schwarz}
    A Kraus map preserves adjoints.  The Kadison--Schwarz inequality therefore
    gives, for every $A\in\MN{D}$,
    \[
        E(A^\dagger A)-E(A^\dagger)E(A)
        =E(A^\dagger A)-E(A)^\dagger E(A)\succeq0.
    \]
\end{proof}

\subsection{Kraus specialization and full algebraic structure}

\begin{definition}[Right and left multiplicative domains]\label{def:one_sided_multiplicative_domains}
    \lean{KadisonSchwarz.rightMultiplicativeDomain,
        KadisonSchwarz.leftMultiplicativeDomain}
    \leanok
    \uses{def:kraus_map}
    For a Kraus map $E$, define
    \[
        \mc{A}_R(E) := \{X \in \MN{D} : \forall Y \in \MN{D},\ E(XY) = E(X) E(Y)\},
    \]
    and
    \[
        \mc{A}_L(E) := \{X \in \MN{D} : \forall Y \in \MN{D},\ E(YX) = E(Y) E(X)\}.
    \]
    We follow the convention that $\mc{A}_R(E)$ controls right multiplication
    and $\mc{A}_L(E)$ controls left multiplication.
    These are Kraus-map specializations of the preceding abstract domains,
    with the names interchanged: here the subscript records the side on which
    the varying factor is appended, whereas the preceding convention records
    the side occupied by the fixed element.
\end{definition}

\begin{definition}[Full multiplicative domain]\label{def:full_multiplicative_domain}
    \lean{KadisonSchwarz.multiplicativeDomain}
    \leanok
    \uses{def:one_sided_multiplicative_domains}
    The \emph{multiplicative domain} of $E$ is
    \[
        \mc{A}(E) := \mc{A}_R(E) \cap \mc{A}_L(E).
    \]
\end{definition}

\begin{remark}[Comparison of the two naming conventions]
    \label{rem:abstract_kraus_multiplicative_domain_names}
    \lean{KadisonSchwarz.mem_abstractRightMultiplicativeDomain_iff,
        KadisonSchwarz.mem_abstractLeftMultiplicativeDomain_iff,
        KadisonSchwarz.mem_abstractMultiplicativeDomain_iff}
    \leanok
    \uses{def:abstract_one_sided_multiplicative_domains,
        def:one_sided_multiplicative_domains,
        def:full_multiplicative_domain}
    Write $\mc{A}_R^{\mathrm W},\mc{A}_L^{\mathrm W}$ for Wolf's
    convention and $\mc{A}_R^{\mathrm K},\mc{A}_L^{\mathrm K}$ for the
    Kraus convention of this subsection.  For the linear map associated with
    a Kraus family,
    \[
        \mc{A}_R^{\mathrm W}=\mc{A}_L^{\mathrm K},\qquad
        \mc{A}_L^{\mathrm W}=\mc{A}_R^{\mathrm K},\qquad
        \mc{A}^{\mathrm W}=\mc{A}^{\mathrm K}.
    \]
\end{remark}

\begin{theorem}[Multiplicative-domain characterization]\label{thm:md_characterization}
    \lean{KadisonSchwarz.mem_rightMultiplicativeDomain_iff, KadisonSchwarz.mem_leftMultiplicativeDomain_iff}
    \leanok
    \uses{def:one_sided_multiplicative_domains, def:unital_kraus}
    Let $E$ be a unital Kraus map. Then
    \[
        X \in \mc{A}_R(E) \iff E(X X^\dagger) = E(X) E(X)^\dagger,
    \]
    and
    \[
        X \in \mc{A}_L(E) \iff E(X^\dagger X) = E(X)^\dagger E(X).
    \]
    Together these are the Kraus-map specialization of
    Theorem~\ref{thm:abstract_md_characterization}.
\end{theorem}

\begin{proof}
    \leanok
    \uses{thm:left_multiplicative_identity}
    If $X$ lies in a one-sided multiplicative domain, evaluate the defining
    identity at $Y = X^\dagger$.
    Conversely, the left-domain implication is exactly
    Theorem~\ref{thm:left_multiplicative_identity}, and the right-domain implication
    follows by applying Theorem~\ref{thm:left_multiplicative_identity} to $X^\dagger$.
\end{proof}

\begin{theorem}[One-sided multiplicative domains are subalgebras]\label{thm:one_sided_md_subalgebras}
    \lean{KadisonSchwarz.rightMultiplicativeDomainSubalgebra, KadisonSchwarz.leftMultiplicativeDomainSubalgebra}
    \leanok
    \uses{def:one_sided_multiplicative_domains, def:unital_kraus}
    For a unital Kraus map $E$, both $\mc{A}_R(E)$ and $\mc{A}_L(E)$ are
    unital subalgebras of $\MN{D}$.
\end{theorem}

\begin{proof}
    \leanok
    \uses{thm:md_characterization}
    Closure under addition follows from linearity of $E$.
    Closure under multiplication: if $X, Y \in \mc{A}_R(E)$, then
    $E((X Y) Z) = E(X) E(Y Z) = E(X) E(Y) E(Z)$ for all $Z$, using
    Theorem~\ref{thm:md_characterization} iteratively, and the
    $\mc{A}_L(E)$ case is analogous.
    Containment of $\Id$ follows from unitality of $E$.
\end{proof}

\begin{theorem}[Multiplicative domain is a $*$-subalgebra]\label{thm:md_star_subalgebra}
    \lean{KadisonSchwarz.multiplicativeDomainStarSubalgebra}
    \leanok
    \uses{def:full_multiplicative_domain, thm:one_sided_md_subalgebras}
    For a unital Kraus map $E$, the full multiplicative domain $\mc{A}(E)$ is
    a $*$-subalgebra of $\MN{D}$.
    This is the algebraic content of \cite[Theorem~5.7]{Wolf2012Quantum}.
\end{theorem}

\begin{proof}
    \leanok
    \uses{thm:md_characterization, thm:one_sided_md_subalgebras}
    If $X \in \mc{A}(E) = \mc{A}_R(E) \cap \mc{A}_L(E)$, then
    $E(X^\dagger X) = E(X)^\dagger E(X)$ and
    $E(X X^\dagger) = E(X) E(X)^\dagger$.
    The characterization of Theorem~\ref{thm:md_characterization} shows
    $X^\dagger \in \mc{A}(E)$.
    Combined with Theorem~\ref{thm:one_sided_md_subalgebras}, this gives a
    $*$-subalgebra.
\end{proof}

\begin{theorem}[Restriction to multiplicative domain is a $*$-homomorphism]\label{thm:md_star_alg_hom}
    \lean{KadisonSchwarz.krausMapStarAlgHom}
    \leanok
    \uses{thm:md_star_subalgebra, def:kraus_map}
    If $E$ is unital, then the restricted map
    $E|_{\mc{A}(E)} : \mc{A}(E) \to \MN{D}$
    is a $*$-algebra homomorphism.
    This is the homomorphism conclusion of \cite[Theorem~5.7]{Wolf2012Quantum}.
\end{theorem}

\begin{proof}
    \leanok
    \uses{thm:kraus_commute}
    Multiplicativity on $\mc{A}(E)$ follows from the definition.
    For $X \in \mc{A}(E)$, the Kraus commutation relation from
    Theorem~\ref{thm:kraus_commute} gives $E(X^\dagger) = E(X)^\dagger$,
    so the restriction preserves adjoints.
\end{proof}


\subsection{Positive maps preserve order and spectral intervals}

\begin{theorem}[Positive maps are monotone]\label{thm:positive_map_monotone}
    \lean{IsPositiveMap.map_le_map}
    \leanok
    \uses{def:positive_map}
    If $T : \MN{n} \to \MN{m}$ is positive and
    $A,B\in\MN{n}$ satisfy $A \le B$, then $T(A) \le T(B)$ in $\MN{m}$.
\end{theorem}

\begin{proof}
    \leanok
    Since $B - A \ge 0$, positivity gives $T(B - A) \ge 0$.
    By linearity, this is $T(B) - T(A) \ge 0$.
\end{proof}

\begin{theorem}[Positive maps preserve adjoints]\label{thm:positive_map_conjTranspose}
    \lean{IsPositiveMap.map_conjTranspose}
    \leanok
    \uses{def:positive_map_is_positive_linear_map}
    If $T : \MN{n} \to \MN{m}$ is positive, then
    $T(A^\dagger) = T(A)^\dagger$ for all $A \in \MN{n}$.
\end{theorem}

\begin{proof}
    \leanok
    Regard $T$ as the associated positive linear map
    (Definition~\ref{def:positive_map_is_positive_linear_map}). For positive
    linear maps between matrix $C^*$-algebras, $T(A^\ast)=T(A)^\ast$.
    Since the star on each matrix algebra is the conjugate transpose,
    $T(A^\dagger)=T(A^\ast)=T(A)^\ast=T(A)^\dagger$.
\end{proof}

\begin{theorem}[Order intervals are preserved]\label{thm:positive_map_image_bounded}
    \lean{IsPositiveMap.image_bounded}
    \leanok
    \uses{thm:positive_map_monotone, def:positive_map}
    Let $T : \MN{D} \to \MN{D}$ be positive with $T(\Id) \le \Id$.
    If $a \le 0 \le b$ and
    $a \Id \le A \le b \Id$,
    then
    \[
        a \Id \le T(A) \le b \Id.
    \]
    This is the matrix form of \cite[Equation~(5.21)]{Wolf2012Quantum}.
\end{theorem}

\begin{proof}\leanok\uses{thm:positive_map_monotone}
    By monotonicity,
    $T(a \Id) \le T(A) \le T(b \Id)$.
    For the lower bound, $T(a \Id) = a T(\Id)$.
    Since $a \le 0$ and $T(\Id) \le \Id$, we have $a T(\Id) \ge a \Id$.
    For the upper bound, $T(b \Id) = b T(\Id) \le b \Id$ since $b \ge 0$.
    Therefore $a \Id \le T(A) \le b \Id$.
\end{proof}

\begin{theorem}[Spectral interval contractivity]\label{thm:positive_map_spectrum_contractivity}
    \lean{IsPositiveMap.spectrum_contractivity}
    \leanok
    \uses{thm:positive_map_image_bounded}
    Let $T : \MN{D} \to \MN{D}$ be positive with $T(\Id) \le \Id$, and let
    $A \in \MN{D}$ be Hermitian.
    If $\spec(A) \subseteq [a,b]$ with $a \le 0 \le b$, then
    \[
        \spec(T(A)) \subseteq [a,b].
    \]
    This is \cite[Equation~(5.21)]{Wolf2012Quantum}.
\end{theorem}

\begin{proof}\leanok\uses{thm:positive_map_image_bounded}
    For Hermitian matrices, the inclusion
    $\spec(A) \subseteq [a,b]$ is equivalent to the order bounds
    $a \Id \le A \le b \Id$.
    Apply Theorem~\ref{thm:positive_map_image_bounded} and translate the
    resulting inequalities back into spectral bounds.
\end{proof}
